%% ============================================================
\clearpage
\section{Judge Prompts}
\label{app:prompts}
%% ============================================================

All LLM-as-Judge and LALM-as-Judge metrics use structured prompts with explicit rating rubrics. Below are the judge prompts for Faithfulness \ref{app:prompt-faithfulness}, Conciseness \ref{app:prompt-conciseness}, Conversation Progression \ref{app:prompt-conv-prog}, Speech Fidelity \ref{app:prompt-speech-fidelity} and User Speech Fidelity \ref{app:prompt-user-speech-fidelity}, User Behavioral Fidelity \ref{app:prompt-user-behavior}, Speakability \ref{app:prompt-speakability}, and Transcription Accuracy (Key Entities) \ref{app:prompt-transcription-accuracy}.

\subsection{Shared Prompt Variables}
\label{app:prompt-shared-variables}

The judge prompts in this appendix use placeholders of the form \texttt{\{variable\_name\}} that are substituted at evaluation time. Most placeholders --- \texttt{conversation\_trace}, \texttt{conversation\_turns}, \texttt{intended\_turns\_formatted}, \texttt{tool\_params}, \texttt{agent\_instructions}, \texttt{agent\_role}, \texttt{available\_tools}, \texttt{user\_simulator\_instructions}, etc.\ --- are derived directly from the dataset (per-record agent and user-simulator configuration) or from the recorded conversation. Their construction from the raw event streams is documented in Appendix~\ref{app:log-processing}; in particular, the linearised \texttt{conversation\_trace} and the per-turn \texttt{intended\_*\_turns} / \texttt{transcribed\_*\_turns} fields are produced by the log-merging procedure described there, and their exact source depends on the pipeline type (cascade, hybrid, or S2S).

A small number of placeholders, however, are \emph{not} record-specific: they are shared text fragments injected into multiple prompts so that every judge sees the same explanation of conventions used in the trace. There are three such fragments. \texttt{interruption\_tags\_reference} is identical across pipelines, while \texttt{user\_turns\_disclaimer} and \texttt{assistant\_turns\_disclaimer} each have a cascade variant and an S2S variant; the appropriate variant is selected at runtime based on the pipeline type of the run being evaluated. The full text of each fragment is reproduced verbatim below.

\subsubsection{Interruption tags reference}
\label{app:prompt-shared-interruption-tags}

The string \texttt{interruption\_tags\_reference} is a non-spoken-tag glossary appended to every prompt that consumes a transcript or trace. It documents the inline annotations produced by the log-merging step (Appendix~\ref{app:log-processing}) so that judges do not penalise the assistant for content that was annotated as interrupted or cut off.

\begin{tcolorbox}[promptbox, title={\small \texttt{interruption\_tags\_reference}}]
These are non-spoken metadata tags inserted during post-processing to annotate speech overlap events. They are NOT part of the spoken text.\newline

Tag definitions:
\begin{itemize}
    \item \texttt{[assistant interrupts]} --- The assistant started speaking while the user was still talking. As a prefix on assistant text, it marks the start of overlapping assistant speech. As an inline marker in user text, it marks approximately where in the user's speech the assistant cut in.
    \item \texttt{[user interrupts]} --- The user started speaking while the assistant was still talking. As a prefix on user text, it marks the start of overlapping user speech. As an inline marker in assistant text, it marks approximately where the user cut in.
    \item \texttt{[likely cut off by user]} --- Appears in assistant text. The assistant's speech was probably cut off by the user starting to speak. Text before this tag may not have been fully spoken. Text after this tag was most likely said (the assistant resumed after the interruption).
    \item \texttt{[likely cut off by assistant]} --- Appears in user text. The user's speech was probably cut off by the assistant starting to speak. Text before this tag may not have been fully spoken.
    \item \texttt{[speaker likely cut itself off]} --- The speaker likely stopped on its own, possibly after detecting overlap or for other reasons, then resumed. Text before this tag may not have been fully spoken. Text after is what the speaker said after resuming.
    \item \texttt{[likely interruption]} --- Catch-all for unexplained breaks in assistant speech that could not be attributed to a specific interruption type.
\end{itemize}
\end{tcolorbox}

\subsubsection{User turns disclaimer}
\label{app:prompt-shared-user-disclaimer}

The string \texttt{user\_turns\_disclaimer} clarifies what the ``user'' rows of the trace actually represent for the pipeline being evaluated, and what the assistant is therefore accountable for. Two variants exist; the cascade variant is used for cascade pipelines, and the S2S variant is used for both S2S and hybrid pipelines (i.e.\ any pipeline where the assistant consumes user audio directly rather than an STT transcript).

\begin{tcolorbox}[promptbox, title={\small \texttt{user\_turns\_disclaimer} (cascade)}]
\textbf{About user turns:} User turns are \textbf{transcripts} produced by the assistant's speech-to-text (STT) system. The assistant receives these transcripts as text input --- this is the only representation of user speech available to the assistant. STT transcripts may contain errors (misheard words, garbled names, dropped syllables), but the assistant cannot know what the user actually said beyond what the transcript shows. Evaluate the assistant against the transcript: if the transcript says ``Kim'' (even if the user actually said ``Kin''), the assistant is acting on ``Kim'' --- that is what it received. Do not penalize the assistant for the transcript's accuracy.
\end{tcolorbox}

\begin{tcolorbox}[promptbox, title={\small \texttt{user\_turns\_disclaimer} (S2S / hybrid)}]
\textbf{About user turns:} This is a \textbf{speech-to-speech} system --- the assistant receives raw audio directly, not a text transcript. The user turns shown here are the \textbf{intended text} (what the user simulator was instructed to say), not what the assistant heard. The assistant is responsible for its own audio understanding. If the assistant misheard the user and acted on incorrect information, that reflects on the assistant --- accurate audio understanding is part of its responsibility. The only mitigation is proper disambiguation: if the assistant was unsure about what it heard, it should have asked the user to confirm or clarify.
\end{tcolorbox}

\subsubsection{Assistant turns disclaimer}
\label{app:prompt-shared-assistant-disclaimer}

The string \texttt{assistant\_turns\_disclaimer} plays the symmetric role for the ``assistant'' rows of the trace, telling the judge whether those rows are the LLM's intended text or an STT transcription of the assistant's audio. As above, two variants exist and pipeline-based dispatch is automatic.

\begin{tcolorbox}[promptbox, title={\small \texttt{assistant\_turns\_disclaimer} (cascade)}]
\textbf{About assistant turns:} Assistant turns shown here are the LLM's \textbf{intended text} --- exactly what the agent produced before TTS rendering. When a user response in the transcript appears to dispute, contradict, or react oddly to an assistant turn that itself looks correct, the most likely cause is an STT error on the user side (the user actually heard something different from what the transcript shows the assistant said). Do not penalize the assistant's prior question, statement, or read-back as ``confusing'' or ``poorly phrased'' in that case --- the assistant LLM had no way to know what the user actually said or heard beyond the transcript.
\end{tcolorbox}

\begin{tcolorbox}[promptbox, title={\small \texttt{assistant\_turns\_disclaimer} (S2S)}]
\textbf{About assistant turns:} This is a \textbf{speech-to-speech} system --- the agent produces audio directly, with no separate intended-text step. The assistant turns shown here are \textbf{STT transcriptions of the agent's audio}, not text the LLM wrote. Audio articulation fidelity (whether the agent \emph{spoke} an entity clearly and correctly) is scored separately by the \texttt{agent\_speech\_fidelity} metric on the actual audio --- do not penalize the agent here for what may be TTS-rendering or STT-transcription artifacts in its turns. Tool call parameters and tool responses shown in the trace are the literal values the agent sent and received via the API, not audio --- if a tool parameter looks wrong, the agent really sent it that way; if the agent's claim contradicts a tool response, the tool truly returned the value shown.
\end{tcolorbox}

\subsection{Faithfulness}
\label{app:prompt-faithfulness}

\begin{tcolorbox}[promptbox, title={\small Faithfulness Prompt}]

      You are an expert evaluator analyzing whether a voice assistant remains faithful to information, policies, and instructions throughout a conversation. You will evaluate the conversation across five dimensions, each scored as a binary flag (true = issue present, false = no issue) and, when flagged, a severity rating (1 = major user impact, 2 = minor or ambiguous; 3 = no issue). See the Rating Scale section below for detailed severity guidance.\newline

      Each dimension evaluates a \textbf{different type of faithfulness violation}. Every issue in the conversation maps to exactly one dimension --- there is no overlap.\newline

      \#\# Agent Instructions\newline
      \{agent\_instructions\}\newline

      \#\# Agent Role\newline
      \{agent\_role\}\newline

      \#\# Available Tools\newline
      \{available\_tools\}\newline

      \#\# Current Date/Time\newline
      \{current\_date\_time\}\newline

      \#\# Understanding User and Assistant Turns\newline
      \{user\_turns\_disclaimer\}\newline

      \{assistant\_turns\_disclaimer\}\newline

      \#\# Full Conversation with Tools\newline
      \{conversation\_trace\}\newline

      \textbf{IMPORTANT --- Interruption tags:} The transcript may contain inline tags indicating speech overlap. These are informational metadata about the voice interaction --- they are NOT faithfulness issues by themselves:\newline

      \{interruption\_tags\_reference\}\newline

      When evaluating, treat these tags as natural voice interaction phenomena. Do NOT treat interrupted or truncated speech as a faithfulness violation. Specifically:
      \begin{itemize}
        \item Do NOT flag truncated content caused by interruptions as hallucination or misrepresentation --- the assistant may not have been able to finish what it was saying.
        \item Do NOT flag incomplete information delivery caused by the user interrupting mid-sentence as a policy violation --- the assistant was cut off, not choosing to withhold information.
        \item Only flag a faithfulness issue if the assistant's actual chosen content (not content lost to interruptions) violates one of the evaluation dimensions.
      \end{itemize}

      \#\# Evaluation Dimensions\newline

      \#\#\# 1. fabricating\_tool\_parameters\newline
      \textbf{Scope: Tool call inputs only.} Did the assistant make a tool call with parameters that were not grounded in user-provided information or prior tool results?\newline

      \textbf{IS a flag:}
      \begin{itemize}
        \item Using a confirmation number, ID, or value that the user did not provide and no prior tool returned
        \item Guessing or inventing parameter values instead of asking the user --- including fabricated IDs and placeholder values like ``?'', ``UNKNOWN'', ``MISSING'', or ``N/A''
        \item Using a parameter value from a different context or conversation segment where it does not apply
        \item Incorrectly categorizing data for enum/categorical tool parameters when the categorization is not supported by the data (e.g., choosing an ``over-threshold'' bucket when the actual value falls below the threshold; or choosing a higher severity/priority category than the described situation supports)
        \item A tool call parameter that cannot be traced to any user statement or prior tool result is a fabrication --- even if the tool happens to return correct results
        \item Hallucinated details in free-text tool fields (e.g., issue\_summary, transfer notes) that were not provided by the user or returned by any tool
        \item Adding random characters to a confirmation number or doubling arbitrary characters to get to the right number of characters.
      \end{itemize}

      \textbf{Is NOT a flag:}
      \begin{itemize}
        \item Using parameter values explicitly stated by the user
        \item Using parameter values returned by a prior tool response (e.g., reusing an ID or record returned by an earlier lookup in a subsequent call)
        \item Using reasonable defaults that are standard for the tool (e.g., a date format conversion)
        \item Standard domain mappings from user-stated information (e.g., ``Chicago O'Hare'' $\rightarrow$ ``ORD'', ``Miami'' $\rightarrow$ ``MIA''; or other unambiguous geographic, enterprise, or industry-standard mappings present in the agent's domain) --- unambiguous mappings are considered grounded
        \item Parameters grounded in policy entitlements derived from prior tool results (e.g., setting an entitlement/waiver/priority flag to the value that the user's status or situation qualifies for per policy, when that qualification is clearly supported by a prior tool result)
        \item Reasonable contextual inferences for categorical parameters, when the enum value's meaning is clearly supported by the user's stated intent
        \item Numeric values derived from prior tool results through simple arithmetic (e.g., summing line items, subtracting used amounts from a balance, or computing a total from prior results)
        \item System-level or framework-generated tool calls made before the assistant has any user input, if the assistant subsequently asks for proper information
        \item Trying valid variations of a user-stated identifier or name when a prior read-only lookup failed (e.g., swapping commonly confused letters like ``v''/``z'' or ``b''/``d'', reordering characters, trying alternative spellings) --- this is proper disambiguation behavior, not fabrication, since the variation is grounded in the user's original utterance
      \end{itemize}

      \textbf{Before flagging a parameter as fabricated:} Verify it cannot be traced to ANY source --- user statements, prior tool results, policy entitlements, simple arithmetic from known values, or standard domain mappings. Also verify enum values against the actual tool specification before claiming a value is invalid.\newline

      \textbf{Read-only vs write tools --- severity guidance:} A fabricated parameter passed to a \textbf{read-only} tool (search, fetch, status check, eligibility check) that fails harmlessly, is self-corrected, and never surfaces incorrect information to the user is still a flag, but rated \textbf{2} (see Rating Scale). A fabricated parameter passed to a \textbf{write} tool --- one that creates, modifies, cancels, or otherwise changes state --- is rated \textbf{1} regardless of whether the tool happens to succeed, because the wrong value reached the system of record.\newline

      \#\#\# 2. misrepresenting\_tool\_result\newline
      \textbf{Scope: How the assistant reports tool results to the user.} Did the assistant inaccurately convey information that was returned by a tool?\newline

      \textbf{IS a flag:}
      \begin{itemize}
        \item Stating incorrect values for fields that the tool response explicitly provided (e.g., wrong departure time, wrong fare amount, wrong seat number, wrong appointment date, wrong ticket priority, wrong PTO balance)
        \item Contradicting what a tool response returned (e.g., saying a flight is on time when the tool showed a delay, stating ``window seat'' when tools show an aisle seat, telling a user their ticket was resolved when status is still open, or reporting a prescription refill was approved when the tool shows denied)
        \item Omitting critical information from a tool result that changes the meaning (e.g., not mentioning a cancellation fee when the tool returned one and a non-zero amount is being charged; not mentioning a side-effect warning returned by a medication lookup; not mentioning that an access request requires manager approval)
        \item Failing to disclose costs/fees or other user-consequential information shown in tool results that the user would need to make an informed decision (e.g., when a non-zero charge is being collected, or when a tool returns eligibility/approval caveats)
        \item Arithmetic errors when computing values from tool data (e.g., incorrectly calculating fare differences, arrival times, remaining balances, or proration amounts) --- verify all math carefully before flagging or clearing
      \end{itemize}

      \textbf{Is NOT a flag:}
      \begin{itemize}
        \item Minor rounding or formatting differences that don't change the meaning (e.g., ``\$384.00'' vs ``\$384'')
        \item Omitting non-essential details from a tool result while accurately conveying the key information
        \item Paraphrasing tool results in conversational language while preserving accuracy
        \item Failing to mention a fee or cost that was \$0 or fully waived (i.e., no amount is being charged), when the financial outcome is accurately communicated
        \item Filtering tool results based on user-stated constraints (e.g., showing only 4 of 5 flights when the 5th doesn't meet the user's arrival time requirement; showing only weekday appointment slots when the user specified weekdays only) --- this is correct behavior, not misrepresentation
        \item Reasonable inferences combining tool data with contextual information (e.g., inferring a flight has departed when scheduled departure is before current time and status shows no cancellation; inferring an SLA breach when a ticket's due date has passed and status is still open)
        \item Time format conversions (e.g., 16:40 = 4:40 PM, 17:00 = 5:00 PM)
      \end{itemize}

      \{misrepresentation\_pipeline\_note\}\newline

      \textbf{Verification requirements:} When checking the assistant's statements against tool results: (1) carefully compute differences/totals (e.g., fare differences as (new fare - original fare) + fees) rather than confusing a total with a delta; (2) check time-format and unit conversions (24h $\leftrightarrow$ 12h, local vs.\ UTC, currency); (3) verify arithmetic independently before flagging a discrepancy; (4) cross-reference ALL relevant tool result fields, not just one.\newline

      \#\#\# 3. violating\_policies\newline
      \textbf{Scope: Agent instructions and policies only.} Did the assistant act in a way that contradicts the agent instructions, system policies, or procedural requirements?\newline

      \textbf{IS a flag:}
      \begin{itemize}
        \item Failing to follow explicit procedural steps outlined in agent instructions (e.g., skipping a required verification step)
        \item Offering options or taking actions that the agent instructions explicitly prohibit
        \item Not applying policies that are clearly applicable to the situation (e.g., not offering an entitled benefit, not following a required disclosure)
        \item Stating a policy incorrectly, or significantly changing a policy's meaning
        \item \textbf{Temporal sequencing for consequential actions:} When instructions require ``explain before acting'' or ``get explicit confirmation before proceeding,'' the assistant must pause for user confirmation BETWEEN read operations and write operations that have financial consequences or are irreversible. Executing such read and write operations in the same turn without intermediate user confirmation violates these instructions. Summarizing results TO the caller after the fact does NOT satisfy a requirement to get confirmation FROM the caller before acting.
        \item \textbf{Irreversible write operations} (cancellations, rebookings, refunds, account/record changes, ticket submissions, access grants, etc.) executed (a) without disclosing what the agent's policies (see Agent Instructions above) require to be surfaced before that write (e.g., a fee, balance impact, eligibility caveat, or approval requirement --- varies by policy), or (b) without a clear signal of user intent for that specific action. Do not invent disclosure requirements the policy does not impose.
      \end{itemize}

      \textbf{Is NOT a flag:}
      \begin{itemize}
        \item Following reasonable interpretations of ambiguous instructions
        \item Minor stylistic deviations from instructions that don't affect the outcome (e.g., slightly different wording for a required disclosure)
        \item Actions not covered by any explicit policy or instruction
        \item Proactive issuance of no-cost benefits the user is clearly entitled to (e.g., goodwill compensation during a service disruption, waiving a fee the user is explicitly entitled to) without explicit confirmation --- these are beneficial actions with no negative consequence, and the user's entitlement or explicit request serves as sufficient basis
        \item When a user explicitly requests a specific action AND the general cost structure has been communicated, proceeding without re-stating exact amounts (if not yet knowable) is not a clear violation
        \item Proceeding with a write tool based on implicit intent (the user described the action they want and provided the necessary information, but did not give a verbatim ``yes, do it now''), when the agent's policies do not require explicit confirmation for that action
        \item ``Predict-then-correct'': stating an action should have no fee/cost based on policy \textit{before} calling the tool, then proceeding to call the tool --- no violation \textbf{when the tool result confirms the prediction} (e.g., ``there should be no fee based on policy'' $\rightarrow$ tool call $\rightarrow$ fee is waived as predicted). If the tool result contradicts the prediction (e.g., a fee is charged anyway), that \textit{is} a faithfulness issue: the user was committed to the action under incorrect cost expectations, and an after-the-fact correction doesn't undo it.
      \end{itemize}

      \textbf{Evaluating policy application:} When two policy paths could apply (e.g., same-day change vs.\ missed flight; in-warranty repair vs.\ out-of-warranty replacement), consider timeline and eligibility carefully. If the user is within the more favorable policy window (e.g., a flight hasn't departed yet, a request is within SLA, a return is within the return window), applying the more favorable applicable policy is not a violation. Also, if two policy paths produce the identical fee/outcome, choosing one over the other is not a material violation.\newline

      \#\#\# 4. failing\_to\_disambiguate\newline
      \textbf{Scope: Handling of ambiguous or contradictory information.} Did the assistant make assumptions or proceed without clarification when the user's input was ambiguous or contradictory? \{disambiguation\_context\}\newline

      \textbf{IS a flag:}
      \begin{itemize}
        \item Proceeding with an action when the user's request could reasonably refer to multiple options and the assistant did not ask which one
        \item Making assumptions about user intent when the user provided contradictory information (e.g., user says two different dates)
        \item Choosing between conflicting pieces of information without asking the user to clarify
        \item Not clarifying ambiguous input that has an impact on the downstream conversation. For example, ``after noon'' and ``afternoon'' could refer to different times of day and should not be silently inferred. The agent should not make a decision that excludes available options without validating the user's intent.
        \item When unable to retrieve some information, if the conversation contains multiple differing versions of a confirmation code or name, the assistant should actively disambiguate rather than silently defaulting to one version or the latest one. Making look-up tool calls is inexpensive and should be done to resolve any ambiguity.
        \item Failing to consider possible input errors when a lookup fails for an uncommon name or alphanumeric code (e.g., not asking the user to spell it out or verify)
        \item Not leveraging required information, such as specific confirmation number or names, that could be reasonably inferred from the conversation.
      \end{itemize}

      \textbf{Is NOT a flag:}
      \begin{itemize}
        \item Proceeding when the user's intent is clear and unambiguous
        \item Asking a clarifying question when the user's request is ambiguous (this is correct behavior)
        \item Making a reasonable inference when the context makes the intent obvious (e.g., user says ``my flight'' when they only have one flight)
        \item Retrying a lookup with a corrected spelling after the user confirms or spells out the information --- this is proper disambiguation behavior
        \item Trying valid different combinations of names and confirmation codes when a lookup fails (e.g., swapping commonly confused letters like ``v''/``z'' or ``b''/``d'', reordering characters)
      \end{itemize}

      \#\#\# 5. hallucination\newline
      \textbf{Scope: Information the assistant states to the user that has no source --- not already covered by the preceding dimensions.} Did the assistant present information that was not provided by the user, not returned by any tool response, and not stated in the agent instructions or system context?\newline

      \textbf{IS a flag:}
      \begin{itemize}
        \item Stating facts, details, or numbers that do not appear in any tool response, user utterance, agent instruction, or system context (e.g., inventing a gate number, adding a benefit the user doesn't have, fabricating an ID or reference number, inventing an amount or deadline)
        \item Presenting fabricated policies, timelines, or conditions not found in any available source
        \item Claiming the system can perform lookups or actions using identifiers not supported by any available tool (e.g., offering to look up a record by an identifier the tools don't accept, or offering a capability not present in the available tools)
        \item Misidentifying the brand, company, or agent role (e.g., using a different airline name, naming the wrong organization)
      \end{itemize}

      \textbf{Is NOT a flag:}
      \begin{itemize}
        \item Stating information that is directly inferable from tool results and/or system context (e.g., computing an arrival time from departure + duration, calculating an expiration date from current date + valid\_months, or computing a remaining amount from a limit minus used amount)
        \item Referencing the current date/time from the system context --- this is grounded information, NOT hallucination
        \item Providing general conversational courtesies that don't assert factual claims
        \item Hedged, commonsense caveats (e.g., ``you may want to verify at the counter'') that don't contradict tool results or policy --- only flag fabricated information presented as definitive fact
        \item General domain knowledge (e.g., standard check-in windows, typical appointment lead times, standard password-reset flows) that is reasonable and not contradicted by tool results
      \end{itemize}

      \textbf{Critical verification step:} Before flagging hallucination, check ALL available sources: (1) all tool responses in the conversation, (2) user utterances, (3) agent instructions, (4) the Current Date/Time field and other system context metadata --- do NOT assume these fields are empty without verifying. Information derived from system context (e.g., current date) is grounded, not hallucinated.\newline

      \textbf{Disambiguation from other dimensions:}
      \begin{itemize}
        \item If the assistant misquotes, distorts or embellishes a tool result $\rightarrow$ flag under misrepresenting\_tool\_result (the source exists but was reported incorrectly)
        \item If an unsupported capability is offered in passing $\rightarrow$ flag here; if actually attempted via fabricated tool call $\rightarrow$ flag under fabricating\_tool\_parameters
        \item If the assistant states something with NO source at all $\rightarrow$ flag here
      \end{itemize}

      You will focus only on the above dimensions. You will NOT consider conversation flow, task completion, or other criteria outside of faithfulness.\newline

      \#\# Rating Scale\newline
      For all five dimensions, determine if there is evidence that one or more issues should be flagged and rate that dimension based on the following guidelines. Severity hinges on the \textbf{impact on the user} --- both inside the conversation and beyond it.

      \begin{itemize}
        \item \textbf{3} (No faithfulness issues):
        \begin{itemize}
          \item No issue with this dimension
        \end{itemize}

        \item \textbf{2} (Minor or ambiguous faithfulness issues --- low user impact):
        \begin{itemize}
          \item A single isolated issue that does not materially affect the outcome, such as:
          \begin{itemize}
            \item Calling a \textbf{read-only} tool (e.g., a lookup, search, or status tool) with wrong parameters, when the error is caught quickly and no incorrect information reaches the user
            \item A small hallucination or misstatement with no consequence inside the conversation and no downstream effect (e.g., a minor phrasing embellishment that does not alter any decision or action)
            \item A minor ungrounded tool parameter that doesn't affect results
            \item A minor policy deviation that doesn't affect the user's decision-making or their understanding of the policy
            \item Skipping a read-back-and-confirm step on values (dates, names, identifiers) the user stated in preceding turns, when the values written to the tool call match what appears in the user's transcribed utterance. \textbf{This carve-out does NOT apply when the transcript shows signals of a likely transcription/STT divergence} --- e.g., the user explicitly said an earlier readback was wrong, OR the transcript contains two or more different versions of the same value across turns. In those cases the assistant had clear signal that a readback was needed to surface the discrepancy, and skipping it likely committed a wrong value to the system --- rate \textbf{1}, not 2.
          \end{itemize}
          \item Minor instruction-following deviations that do not materially affect the outcome (e.g., slight formatting differences, omitting low-importance optional steps)
          \item Borderline cases where it is unclear whether a faithfulness violation occurred due to ambiguous instructions, incomplete context, or reasonable interpretation differences
          \item Adopting incorrect terminology from the user (e.g., wrong brand or product name) while processing the correct record, when it doesn't cause confusion or incorrect actions
          \item If something appears as being borderline an issue, it should probably be rated 2.
        \end{itemize}

        \item \textbf{1} (Clear faithfulness violations --- major user impact):
        \begin{itemize}
          \item Any issue that materially affects the user, either during the conversation or afterward. Such as:
          \begin{itemize}
            \item Calling a \textbf{write} tool with wrong or fabricated parameters --- regardless of whether the call succeeds.
            \item Committing a value to a write tool when the transcript shows clear signals of a likely transcription/STT divergence (e.g., the user explicitly said an earlier readback was wrong, OR the transcript contains two or more different versions of the same value across turns) and the assistant did not resolve the ambiguity before committing --- \textbf{rate 1 even if the written value happens to match the most recent transcribed utterance}, because the divergence signal made a readback necessary and the assistant likely committed a wrong value.
            \item Executing an \textbf{irreversible action without the explicit user confirmation that the agent's policies (see Agent Instructions) require} for that specific action (e.g., a ``summarize and confirm before submitting'' rule). If the agent's policy does not require explicit confirmation for that action, implicit intent from the conversation flow is sufficient --- do not flag the absence of a verbatim ``yes, do it now.''
            \item Executing an \textbf{irreversible action without first disclosing what the agent's policies (see Agent Instructions) require to be surfaced} before that write (e.g., balance impact, fee, eligibility caveat, approval requirement --- varies by policy). Do not invent disclosure requirements the policy does not impose.
            \item \textbf{Financial or reputational impact} on the user or the company --- e.g., communicating incorrect charges, fees, refunds, balances, or amounts; making commitments the company cannot honor; misstating obligations in a way that could harm trust
            \item Leaving the user with an \textbf{incorrect understanding of company policy} that they could act on in the future (e.g., misstating eligibility rules, entitlement conditions, or process requirements), even if the current conversation ends up fine
            \item Hallucinating information not present in tool results, especially consequential figures (costs, balances, dates, approvals) communicated to the user as fact
          \end{itemize}
          \item Any faithfulness issue that repeatedly prevents the conversation from progressing is also rated 1.
        \end{itemize}
      \end{itemize}

      For the final rating of the conversation, use the minimum rating across all dimensions as the overall faithfulness rating (i.e., if any dimension is rated 1, overall rating is 1; if all dimensions are 3, overall rating is 3; if there are no 1s but at least one 2, overall rating is 2).\newline

      \#\# Response Format\newline
      Respond in JSON format:\newline
      \{\{\newline
          ``dimensions'': \{\{\newline
              ``fabricating\_tool\_parameters'': \{\{\newline
                  ``evidence'': ``\textless{}string\textgreater{}'',\newline
                  ``flagged'': \textless{}bool: true if issue is present, false otherwise\textgreater{},\newline
                  ``rating'': \textless{}int: 1, 2, or 3\textgreater{}\newline
              \}\},\newline
              ``misrepresenting\_tool\_result'': \{\{\newline
                  ``evidence'': ``\textless{}string\textgreater{}'',\newline
                  ``flagged'': \textless{}bool: true if issue is present, false otherwise\textgreater{},\newline
                  ``rating'': \textless{}int: 1, 2, or 3\textgreater{}\newline
              \}\},\newline
              ``violating\_policies'': \{\{\newline
                  ``evidence'': ``\textless{}string\textgreater{}'',\newline
                  ``flagged'': \textless{}bool: true if issue is present, false otherwise\textgreater{},\newline
                  ``rating'': \textless{}int: 1, 2, or 3\textgreater{}\newline
              \}\},\newline
              ``failing\_to\_disambiguate'': \{\{\newline
                  ``evidence'': ``\textless{}string\textgreater{}'',\newline
                  ``flagged'': \textless{}bool: true if issue is present, false otherwise\textgreater{},\newline
                  ``rating'': \textless{}int: 1, 2, or 3\textgreater{}\newline
              \}\},\newline
              ``hallucination'': \{\{\newline
                  ``evidence'': ``\textless{}string: 1--2 sentences citing specific examples from the transcript, or `None' if not flagged\textgreater{}'',\newline
                  ``flagged'': \textless{}bool: true if issue is present, false otherwise\textgreater{},\newline
                  ``rating'': \textless{}int: 1, 2, or 3\textgreater{}\newline
              \}\}\newline
          \}\},\newline
          ``rating'': \textless{}int: 1, 2, or 3 --- minimum rating across all dimensions\textgreater{}\newline
      \}\}
\end{tcolorbox}

\begin{tcolorbox}[promptbox, title={\small \texttt{disambiguation\_context} (Cascade)}]
Since the assistant is working from a speech-to-text transcript, it should account for potential transcription errors, and clarify any ambiguity in the user's intent, especially when they lead to write/irreversible operations. It's not needed to clarify if the tools called are simple lookups, but if the lookups fail, the assistant is expected to clarify the user's intent.
\end{tcolorbox}

\begin{tcolorbox}[promptbox, title={\small \texttt{disambiguation\_context} (S2S / hybrid)}]
Since the assistant processes raw audio directly (speech-to-speech), it should account for potential audio perception errors — mishearing letters, numbers, names, or codes is common with spoken input. The assistant should clarify any ambiguity, especially for alphanumeric codes, names, and values that lead to write/irreversible operations. It's not needed to clarify if the tools called are simple lookups, but if the lookups fail, the assistant is expected to clarify the user's intent. The bar for disambiguation is higher than for a text-based system because the assistant knows it is working from audio and should anticipate mishearings.
\end{tcolorbox}

\begin{tcolorbox}[promptbox, title={\small \texttt{misrepresentation\_pipeline\_note} (S2S)}]
**Speech-to-speech scoping for this dimension.** Because assistant turns in the trace are STT-transcribed audio (see *About assistant turns* above), token-level discrepancies between an assistant utterance and a tool result — dropped/added dashes, single-character substitutions, missing or extra digits within long alphanumeric IDs, altered spacing — typically reflect TTS-rendering or STT-transcription artifacts and are scored by `agent\_speech\_fidelity`, not here. Only flag `misrepresenting\_tool\_result` when the discrepancy is structural/semantic (wrong field, wrong order of magnitude, wrong category) or when downstream signals — subsequent tool calls, 
follow-up actions, user objections — show the agent was internally operating on a wrong value.
\end{tcolorbox}


\subsection{Speech Fidelity}
\label{app:prompt-speech-fidelity}

\begin{tcolorbox}[promptbox, title={\small Speech Fidelity Prompt}]

        You are an expert evaluator judging the fidelity of this audio file against the intended text.
        You will listen to one audio clip and verify that the spoken content faithfully reproduces the intended text, with special attention to TTS-critical entities.
        The audio provided is a recording of the agent's side of a conversation, and contains only the agent responses, not the user.\newline

        \#\# Intended Turns\newline
        \{intended\_turns\_formatted\}\newline

        \#\# IMPORTANT: Comparison Rules\newline

        Your task is to compare the \textbf{exact intended text} word-for-word against what you hear in the audio. The TTS-critical entities highlight which parts are most important to verify, but they do NOT replace or override the intended text.\newline

        \#\# Understanding the Intended Text\newline

        The intended text may contain non-spoken tags and markers. You must understand these to evaluate fairly.\newline

        \#\#\# Audio-Direction Tags\newline
        Tags like [slow], [firm], [annoyed] describe how the words were meant to be spoken. They are NOT spoken aloud and should never be expected in the audio.\newline

        \#\#\# Interruption Tags\newline
        \{interruption\_tags\_reference\}\newline

        The tags tell you that certain portions of the intended text were likely never spoken, because the speaker was interrupted or cut themselves off. Do NOT penalize for missing words that fall in a region the tags indicate was not spoken.\newline

        \textbf{Key principle:} If a tag indicates that a section of text was likely not spoken aloud (due to interruption or cut-off), do NOT penalize for those words being missing from the audio. Only evaluate fidelity for words that were reasonably expected to have been spoken.\newline

        \#\# Evaluation Criteria\newline

        For each intended turn, compare what you hear in the audio against the intended text. Focus especially on \textbf{TTS-critical entities} listed for each turn.\newline

        \textbf{Entity categories to watch:}
        \begin{itemize}
          \item Confirmation codes (e.g., ZK3FFW, FAR0UM, 8JVSDF)
          \item Domain-specific identifiers (e.g., flight numbers like ``SkyWay 410'', ticket or incident numbers, order numbers, case IDs)
          \item Dollar amounts (e.g., \$15, \$1,285.00)
          \item Short alphanumeric codes (e.g., seat numbers like ``21C'', room numbers, extension numbers)
          \item Spelled-out codes (e.g., ``Z K three F F W'') --- verify EVERY letter and digit individually; ``K O L T S F'' vs ``K O L T S S F'' is an error
          \item Reference IDs with segments (e.g., REF-8JVSDF-001, MEAL-FAR0UM-PAX0) --- verify each segment; ``M E L'' vs ``M E A L'' is an error
          \item Times (e.g., 3:55 PM, 10:30 AM)
          \item Dates (e.g., March 25th, February 3rd)
          \item Names (e.g., Mr.\ Rivera, Rodriguez)
        \end{itemize}

        \textbf{What constitutes an error (rating = 0):}
        \begin{itemize}
          \item Any entity spoken incorrectly (wrong digits, letters, amounts, numbers)
          \item Missing words that change the meaning or omit an entity
          \item Added words that introduce a factually incorrect entity
          \item Substituted words that alter an entity value
        \end{itemize}

        \textbf{What to ignore (does NOT cause rating = 0):}
        \begin{itemize}
          \item Minor pronunciation variations that do not change the identity of an entity (e.g., ``Ms.'' vs ``Miss'' is acceptable)
          \item Filler words (``um'', ``uh'', ``so'') added or omitted
          \item End-of-audio cut-off: if the audio cuts off at the very END of the last turn, missing trailing words is acceptable as long as all entities in that turn were spoken correctly before the cut-off
          \item Slight pacing or prosody differences
          \item Non-spoken tags: [slow], [firm], [annoyed], and all interruption tags listed above
          \item Words in regions flagged by interruption tags as likely not spoken
        \end{itemize}

        \#\# Rating Scale (per turn)
        \begin{itemize}
          \item \textbf{1 (High Fidelity)}: All entities are spoken correctly. Non-entity words are faithfully reproduced with no meaningful omissions or additions.
          \item \textbf{0 (Low Fidelity)}: One or more entity errors, OR significant non-entity word errors that change the meaning of the turn.
        \end{itemize}

        \#\# Response Format\newline
        Respond with a JSON object. Each turn entry must include the turn\_id matching the turn number shown in the Intended Turns above:\newline
        \{\{\newline
          ``turns'': [\newline
            \{\{\newline
              ``turn\_id'': \textless{}int: the turn number from the Intended Turns\textgreater{},\newline
              ``transcript'': \textless{}string: your transcription of the audio for this turn, use only the audio for this not the intended text\textgreater{}\newline
              ``explanation'': ``\textless{}string: 1--3 sentence analysis of fidelity for this turn, citing specific intended vs actual mismatches, noting any regions skipped due to interruption flags\textgreater{}'',\newline
              ``rating'': \textless{}0 or 1\textgreater{}\newline
            \}\}\newline
          ]\newline
        \}\}
\end{tcolorbox}

\subsection{Conversation Progression}
\label{app:prompt-conv-prog}

\begin{tcolorbox}[promptbox, title={\small Conversation Progression Prompt}]
      You are an expert evaluator analyzing whether a voice assistant effectively moved a conversation forward. You will evaluate the conversation across four dimensions, each scored as a binary flag (true = issue present, false = no issue).\newline

      Each dimension evaluates a \textbf{different type of action}. Every issue in the conversation maps to exactly one dimension --- there is no overlap.
      Ensure to consider both the assistant agent instructions and the following agent dimensions when evaluating the conversation.\newline

      \textbf{IMPORTANT --- Scope boundary with faithfulness:} This metric evaluates whether the conversation moved forward efficiently. It does NOT evaluate whether the assistant followed policies, complied with user constraints, or acted faithfully to its instructions --- those are faithfulness concerns.
      If an issue is primarily about the assistant violating a policy or acting against the user's explicit instructions (e.g., taking an action the user said not to, not disclosing fees), do NOT flag it here even if it also affected conversation flow. Only flag issues where the assistant's conversational choices (questions asked, information repeated, tools called) were themselves inefficient or counterproductive.\newline

      \textbf{IMPORTANT --- Voice conversation context:} This is a voice (spoken) conversation, which means speech recognition errors are common.
      When the assistant repeats a request because the previous attempt was misheard or garbled, this is expected behavior in a voice interface, not a progression issue.\newline

      \textbf{IMPORTANT --- Interruption tags:} The transcript may contain inline tags indicating speech overlap. These are informational metadata about the voice interaction --- they are NOT conversation progression issues by themselves:\newline

      \{interruption\_tags\_reference\}\newline

      When evaluating, treat these tags as natural voice interaction phenomena. Do NOT penalize interruptions themselves. Only flag an issue if the interruption caused observable consequences (e.g., information loss because the agent's cut-off speech contained critical details that were never restated, or unnecessary repetition because the agent repeated already-heard information after being interrupted).\newline

      \#\# Understanding the Conversation Trace\newline
      \{user\_turns\_disclaimer\}\newline

      \{assistant\_turns\_disclaimer\}\newline

      \#\# Full Conversation with Tools\newline
      \{conversation\_trace\}\newline

      \#\# Evaluation Dimensions\newline

      \#\#\# 1. unnecessary\_tool\_calls\newline
      \textbf{Scope: Tool call actions only.} Were any tool calls unjustified --- repeated without reason, made without required information, or made for data already available?\newline

      \textbf{IS a flag:}
      \begin{itemize}
        \item Calling the same tool with the same parameters after a prior successful response (no new user input or error in between)
        \item Calling a tool with empty or missing required parameters, causing a predictable error (e.g., calling a lookup tool with empty strings before asking the user for the required identifier)
        \item Calling a tool when the needed information was already returned by a previous tool response
        \item Calling a tool to verify something a prior tool response already confirmed
      \end{itemize}

      \textbf{Is NOT a flag:}
      \begin{itemize}
        \item Retrying a tool call after a tool error with corrected parameters
        \item Calling the same tool with different parameters (e.g., different IDs or search criteria)
        \item Sequential tool calls that each return new, necessary information (e.g., a record lookup followed by a status check followed by a related-details fetch)
        \item A tool call that fails unexpectedly (the assistant could not have predicted the failure)
        \item Tool calls that are necessary for the task but were executed prematurely (e.g., before the user confirmed) --- premature execution is a faithfulness/policy compliance issue, not a conversation progression issue
        \item Tool calls that follow standard agent instructions (e.g., automatically carrying over related attributes or defaults when taking an action) even if the user did not explicitly request those specific actions
      \end{itemize}

      \textbf{CAVEAT: If the model makes 3 or more unnecessary tool calls, this dimension should be rated 1.}\newline

      \#\#\# 2. information\_loss\newline
      \textbf{Scope: The assistant's memory of established facts.} Did the assistant fail to retain or act on information already established in the conversation --- whether from the user's statements or from prior tool responses?\newline

      This dimension is about the assistant \textbf{forgetting or ignoring known facts}, regardless of how that failure manifests (re-asking, wrong assumptions, ignoring constraints).\newline

      \{information\_loss\_pipeline\_note\}\newline

      \textbf{IS a flag:}
      \begin{itemize}
        \item Re-asking the user for information they already provided (e.g., asking for the confirmation number or reference ID after the user stated it and it was used successfully).
        \item Ignoring a constraint the user explicitly stated (e.g., the user ruled out a particular action but the assistant still asks about it or asks for the details that would only be needed to take that action)
        \item Failing to use relevant data from a prior tool response when it was needed for the next step (e.g., not using an identifier returned by an earlier lookup when making a follow-up tool call that requires it)
      \end{itemize}

      \textbf{Is NOT a flag:}
      \begin{itemize}
        \item Asking for information the user has not yet provided
        \item Asking a clarifying question about genuinely ambiguous information
        \item Asking for authentication or identification details required by the agent instructions (e.g., a confirmation number, reference ID, or user's name) at the start of the conversation
        \item The assistant acting on information that contradicts what the user said, when the contradiction is due to a faithfulness or policy violation --- flag that under faithfulness, not here. Only flag here if the assistant demonstrably forgot or ignored previously established facts within the conversation flow.
      \end{itemize}

      \textbf{Disambiguation from other dimensions:}
      \begin{itemize}
        \item If the assistant re-asks for user-provided info $\rightarrow$ flag here (not redundant\_statements)
        \item If the assistant makes an unnecessary tool call because it forgot a prior result $\rightarrow$ flag under unnecessary\_tool\_calls (the tool action is the observable problem)
        \item If the assistant proceeds with an action that contradicts the user's stated preference (e.g., choosing a different option than the one the user requested) $\rightarrow$ this is a faithfulness violation, not information\_loss. Only flag here if the assistant clearly forgot the user's input, not if it chose to override it.
      \end{itemize}

      \#\#\# 3. redundant\_statements\newline
      \textbf{Scope: The assistant repeating its own previous output.} Did the assistant restate information it had already communicated to the user?\newline

      This dimension ONLY covers the assistant repeating \textbf{its own prior utterances} --- not forgetting user input (that is information\_loss) and not tool call issues (that is unnecessary\_tool\_calls).\newline

      \textbf{IS a flag:}
      \begin{itemize}
        \item Restating details, times, amounts, or status information the assistant already told the user in an earlier turn (outside of a final recap) when the user did not ask for it
        \item Repeating the same explanation or instruction in multiple turns when the user has acknowledged and moved on
      \end{itemize}

      \textbf{Is NOT a flag:}
      \begin{itemize}
        \item A brief recap or summary at the end of the conversation (this is helpful, not redundant). However, if the assistant provides multiple recaps across different turns, only the final one is exempt --- earlier recaps that restate already-communicated information are still flagged.
        \item Confirming back details to the user once for verification (e.g., reading back a confirmation number the user just provided)
        \item Stating information for the first time, even if it was available from a tool response earlier
        \item Repeating information in direct response to the user explicitly requesting confirmation or asking to hear it again (the user must clearly ask --- simply continuing the conversation is not a request for repetition)
        \item Re-explaining a policy or constraint when the user continues to challenge, dispute, or insist against it --- the assistant must reiterate its position in these cases and should not be penalized for doing so. However, if the assistant repeats the exact same explanation verbatim across multiple turns, flag it --- the assistant should vary its phrasing.
        \item Repeating a request for information (e.g., confirmation code, spelling) when speech recognition or transcription errors clearly caused the previous attempt to fail (e.g., garbled text, partial characters, obvious mishearing visible in the transcript). Do NOT apply this exception when the transcript shows no evidence of ASR failure --- the assistant re-asking without cause is still a flag.
      \end{itemize}

      \#\#\# 4. question\_quality\newline
      \textbf{Scope: The quality and appropriateness of the assistant's questions, where the issue is NOT caused by forgetting information (that is information\_loss).} Did the assistant ask poorly formed questions or fail to ask a necessary clarifying question?\newline

      \textbf{IS a flag:}
      \begin{itemize}
        \item Asking an overly broad or vague question when the assistant had enough information to take action (e.g., ``What would you like to do?'' when the user already stated a clear goal that the assistant remembers but chose not to act on)
        \item Asking multiple questions at once when a single tool call could have resolved the need
        \item Failing to ask for clarification when the user's request was genuinely ambiguous, and instead proceeding with assumptions
        \item Failing to ask for clarification when there are multiple options that meet the user's requirements
        \item Failing to ask for required information before taking an action (e.g., not asking for required details for a tool call before making the tool call, when those details have not been made available through a previous tool call, or inputs from the user)
        \item Failing to provide necessary information for the user to make a decision (e.g., not providing clear information about the details of the options available to the user)
        \item Taking an irreversible action (e.g., cancellation, rebooking, ticket submission, access grant, account change) without first confirming when user input is ambiguous or contradicts system data (e.g., user claims a 4-hour delay but system shows 45 minutes --- assistant should clarify before acting)
      \end{itemize}

      \textbf{Is NOT a flag:}
      \begin{itemize}
        \item Asking for required authentication or identification information required by the agent instructions (e.g., confirmation number, reference ID, user's name)
        \item Asking a clarifying question when the user's intent is genuinely ambiguous
        \item Asking a follow-up question based on new information from a tool response
        \item Asking the user to confirm an error-prone value (alphanumeric ID, code, date, dollar amount, spelled-out name) just provided --- read-backs of error-prone values are standard voice-agent practice required by typical agent policies and are NOT a question quality issue. This holds even when the transcript shows the assistant's read-back as matching what the user said: STT may render different spoken audio as identical text, so the read-back can still be catching a real audio-level mismatch.
        \item Not disclosing fees, costs, or other policy-required details before taking an action --- policy compliance (e.g., whether the assistant explained consequences before an irreversible action) is a faithfulness concern, not a conversation progression issue. This dimension only evaluates whether the assistant's questions and information-sharing effectively moved the conversation forward.
        \item Referencing information that exists in the agent instructions (e.g., standard fees, policies) without verifying it via a tool call --- the agent is expected to know its own instructions. Only flag if the information was genuinely unknown and required a tool call or user input.
      \end{itemize}

      \textbf{Disambiguation from information\_loss:}
      \begin{itemize}
        \item If the assistant asks ``What would you like to do?'' because it FORGOT the user already stated their goal $\rightarrow$ flag under information\_loss
        \item If the assistant asks ``What would you like to do?'' when the user's goal is clear and remembered but the assistant chose a vague question over taking action $\rightarrow$ flag here
      \end{itemize}

      \#\# Rating Scale\newline
      For all four dimensions, determine if there is evidence that one or more issues should be flagged and rate that dimension based on the following guidelines:

      \begin{itemize}
        \item \textbf{3} (No progression issue):
        \begin{itemize}
          \item No issue with this dimension
        \end{itemize}
        \item \textbf{2} (Minor progression issue):
        \begin{itemize}
          \item A single isolated issue that does not significantly impact the conversation flow (e.g., one unnecessary tool call that didn't slow things down, a single redundant restatement, one vague question)
          \item A borderline case where it is unclear whether the issue constitutes a real progression problem
        \end{itemize}
        \item \textbf{1} (Clear progression issue):
        \begin{itemize}
          \item Multiple instances of the same type of issue in this dimension
          \item A single severe issue that clearly derailed or stalled the conversation (e.g., ignoring a stated constraint or user requirement before carrying out a write operation, failing to ask for required information before taking action, asking an overly vague question when the user's goal was clear, making an overly vague assumption not supported by user inputs/conversation history when multiple options exist)
        \end{itemize}
      \end{itemize}

      \#\# Overall Rating\newline
      The final rating considers BOTH the severity within each dimension AND the total number of flagged dimensions:

      \begin{itemize}
        \item \textbf{3}: No dimension is flagged (all dimensions rated 3)
        \item \textbf{2}: One or two dimensions are flagged at rating 2 (minor), AND no dimension is rated 1
        \item \textbf{1}: Any of the following:
        \begin{itemize}
          \item Any dimension is rated 1 (clear issue within a single dimension)
          \item Three or more dimensions are flagged (even if each is individually minor, widespread issues across many areas constitute a clear overall progression problem)
        \end{itemize}
      \end{itemize}

      \#\# Response Format\newline
      Respond in JSON format. The ``evidence'' field must ALWAYS contain 1--2 sentences referencing specific parts of the transcript, even when flagged is false. When not flagged, briefly explain why no issue was found.\newline
      \{\{\newline
          ``dimensions'': \{\{\newline
              ``unnecessary\_tool\_calls'': \{\{\newline
                 ``evidence'': ``\textless{}string: REQUIRED --- cite transcript examples if flagged, or explain why clean if not\textgreater{}'',\newline
                 ``flagged'': \textless{}bool: true if issue is present, false otherwise\textgreater{},\newline
                 ``rating'': \textless{}int: 1, 2, or 3\textgreater{}\newline
              \}\},\newline
              ``information\_loss'': \{\{\newline
                  ``evidence'': ``\textless{}string: REQUIRED\textgreater{}'',\newline
                  ``flagged'': \textless{}bool: true if issue is present, false otherwise\textgreater{},\newline
                  ``rating'': \textless{}int: 1, 2, or 3\textgreater{}\newline
              \}\},\newline
              ``redundant\_statements'': \{\{\newline
                  ``evidence'': ``\textless{}string: REQUIRED\textgreater{}'',\newline
                  ``flagged'': \textless{}bool: true if issue is present, false otherwise\textgreater{},\newline
                  ``rating'': \textless{}int: 1, 2, or 3\textgreater{}\newline
              \}\},\newline
              ``question\_quality'': \{\{\newline
                  ``evidence'': ``\textless{}string: REQUIRED\textgreater{}'',\newline
                  ``flagged'': \textless{}bool: true if issue is present, false otherwise\textgreater{},\newline
                  ``rating'': \textless{}int: 1, 2, or 3\textgreater{}\newline
              \}\}\newline
          \}\},\newline
          ``rating'': \textless{}int: 1, 2, or 3\textgreater{}\newline
      \}\}
\end{tcolorbox}

\begin{tcolorbox}[promptbox, title={\small \texttt{information\_loss\_pipeline\_note} (S2S)}]
**Speech-to-speech scoping for this dimension.** Because assistant turns in the trace are STT-transcribed audio (see *About assistant turns* above), variant token-level readings of the same alphanumeric identifier across nearby assistant turns — dropped/added dashes, single-character substitutions, missing or extra digits within long IDs, altered spacing or capitalization — typically reflect TTS-rendering or STT-transcription artifacts on a value the agent is reading consistently in audio. These are scored by `agent\_speech\_fidelity`, not here. Only flag `information\_loss` when the discrepancy is structural/semantic (different entity, wrong field, wrong category — e.g., addressing the user by an entirely different first name, or referencing a different person/record than the tool returned), or when downstream signals — subsequent tool calls made with a wrong value, follow-up actions taken on stale data, user objections that the agent then fails to incorporate — show the agent was internally operating on a wrong value or had genuinely lost track of the established fact."
\end{tcolorbox}

\subsection{Conciseness}
\label{app:prompt-conciseness}

\begin{tcolorbox}[promptbox, title={\small Conciseness Prompt}]
      You are an expert evaluator judging the conciseness and voice-appropriateness of assistant responses in a voice conversation.\newline

      \#\# Conversation\newline
      \{conversation\_turns\}\newline

      \#\# Understanding the Conversation Format\newline

      The conversation is grouped by turn\_id. Each turn may contain:
      \begin{itemize}
        \item \textbf{user}: What the user said
        \item \textbf{assistant}: What the assistant said (there may be multiple assistant entries within a single turn --- e.g., the assistant speaks, calls a tool, then speaks again)
        \item \textbf{tool\_call}: A tool invocation made by the assistant
        \item \textbf{tool\_response}: The result returned by the tool
      \end{itemize}

      When a turn contains multiple assistant entries, evaluate them \textbf{together as a single unit} --- they represent the assistant's complete response within that turn. Tool calls and responses between assistant entries explain why the assistant spoke in multiple parts (it was waiting for data). It could also be due to interruptions from the user.\newline

      \#\# Understanding Interruption Tags\newline

      \{interruption\_tags\_reference\}\newline

      \textbf{Key principle:} When interruption tags are present, the assistant may not have been able to finish what it was saying. Do NOT penalize for truncated or fragmented content caused by interruptions. Only evaluate the conciseness of content the assistant chose to say, not content that might have been cut off.\newline

      \#\# Instructions\newline
      The conversation includes user, assistant, tool\_call, and tool\_response entries. Rate only the assistant's spoken content. User turns, tool calls, and tool responses are provided for context only.\newline

      For each turn that contains assistant content, evaluate whether the assistant's response is appropriately concise and easy to digest when spoken aloud to a human.\newline

      The assistant is expected to follow conversational voice guidelines:
      \begin{itemize}
        \item Keep responses brief and conversational (typically 2--4 sentences)
        \item Summarize long lists rather than reading them exhaustively
        \item Avoid overwhelming the listener with too much information at once
        \item Spread multiple requests across turns when possible
        \item Present options conversationally and avoid cramming excessive detail into one turn
      \end{itemize}

      \#\# Evaluation Criteria\newline
      When evaluating each turn, consider:
      \begin{itemize}
        \item Does the response get to the point without filler, rambling, or unnecessary content?
        \item Is all the information relevant and necessary given the conversation context?
        \item Is the amount of detail reasonable for someone listening to --- not reading --- the response?
        \item If the response enumerates options or items (e.g., ``Option one is\ldots{} Option two is\ldots{}''), does the structure help the user? The volume should not be overwhelming.
        \item Is the provided information justified by context (e.g., confirming a detail the user may have misheard)? Or is it inappropriate (e.g., excessive itemization or explanation when the user may only care about the end result)?
        \item Within turns, is repetition avoided? Across turns there may be valid reasons for repetition, but it should usually not occur within a single turn.
        \item Essential information --- such as confirmation codes, reference IDs, ticket numbers, or other specific details the user needs to note down --- should never be penalized, regardless of length.
      \end{itemize}

      \#\# Allowed Exceptions (Voice Interaction Realities)\newline
      The assistant may occasionally produce longer turns when the context requires precise information transfer. The following cases should NOT be penalized for verbosity or information density. The turn itself may still be penalized for other reasons.

      \begin{enumerate}
        \item \textbf{Phonetic Confirmation of Codes}
        \begin{itemize}
          \item When confirming a confirmation code, reference number, ticket number, or similar identifier, the assistant may spell characters using the NATO phonetic alphabet (e.g., ``B as in Bravo, F as in Foxtrot'').
          \item This is especially appropriate when the user previously misheard or asked for clarification.
        \end{itemize}
        \item \textbf{Reference or Identifier Delivery}
        \begin{itemize}
          \item When providing an identifier the user needs to note down (e.g., a ticket number, reference code, or voucher code), the assistant may read the whole code out loud.
          \item This information is essential and should not be penalized regardless of length.
        \end{itemize}
        \item \textbf{End-of-Call Wrap-Up}
        \begin{itemize}
          \item The final assistant turn in a conversation may include a slightly longer recap or confirmation of next steps (e.g., summarizing the action taken, confirming what will be sent or followed up on, thanking the user).
          \item Minor additional detail in this final wrap-up should not be penalized unless it becomes excessively long or introduces unrelated information.
        \end{itemize}
      \end{enumerate}

      Important principle: Information given in assistant turns must be short enough for an average person to easily follow in real-time conversation and retain in working memory.\newline

      \#\# Failure Modes\newline
      When a response is not optimally concise, identify which of the following failure modes are present. A turn may have multiple failure modes.\newline

      \textbf{verbosity\_or\_filler}\newline
      Contains unnecessary wording, repetition within the same turn, hedging, or explanation beyond what the context requires.\newline

      \textbf{excess\_information\_density}\newline
      Presents too many distinct facts, options, numbers, steps, or requests at once, making it difficult for a listener to process in real time. Note: bundling closely related transactional details that the user needs to act on or remember together (e.g., confirming a reference number, date, and one or two key details in a single turn) is expected behavior --- only flag this when the volume of information genuinely exceeds what a listener can comfortably retain.\newline

      \textbf{over\_enumeration\_or\_list\_exhaustion}\newline
      Reads out long lists instead of summarizing, or presents multiple options with excessive detail rather than inviting follow-up.\newline

      \textbf{contextually\_disproportionate\_detail}\newline
      Provides more background, clarification, or explanation than the situation warrants.\newline

      \#\# Contextual Leniency and Failure Mode Priority\newline
      Conciseness should be evaluated with respect to the conversational context. If additional wording or detail is clearly necessary for the user to understand or act on the information, a modest increase in verbosity should be considered acceptable and should NOT be penalized.\newline

      If none of the above are present, return an empty list for failure\_modes.\newline

      \#\# Rating Scale For Each Turn With Assistant Content
      \begin{itemize}
        \item \textbf{3} (Highly Concise / No Cognitive Overload) --- The response is clear, appropriately scoped for voice, and comfortably digestible in real time. No failure modes are present. A turn that delivers a few closely related facts as part of a single transactional step (e.g., confirming the key details of a request or incident) still qualifies as 3 if the listener can comfortably absorb it in one pass.
        \item \textbf{2} (Adequate but Not Optimally Concise) --- One minor failure mode is present, but the response remains reasonably processable in a voice setting and does not meaningfully overwhelm the listener. Reserve this rating for turns where you can identify specific content that should have been omitted or deferred to a later turn --- not merely for turns that happen to contain several necessary details.
        \item \textbf{1} (Not Concise / Causes Cognitive Overload) --- One or more significant failure modes are present that materially increase cognitive load and would hinder comprehension in a voice conversation.
      \end{itemize}

      Provide one entry per turn\_id in the conversation.\newline

      \#\# Response Format\newline
      Provide your response as a valid JSON array, one entry per turn. Each entry must include the turn\_id matching the turn number shown in the conversation above.
      \begin{itemize}
        \item If the turn contains assistant content, rate it with 1, 2, or 3.
        \item If the turn does not contain assistant content (e.g., user-only turn), set rating to null.
      \end{itemize}
      [\newline
        \{\{\newline
          ``turn\_id'': \textless{}int: the turn number from the conversation\textgreater{},\newline
          ``explanation'': ``\textless{}Detailed analysis referencing the evaluation criteria and explicitly linking identified weaknesses to the listed failure modes to justify the selected rating (1--3). Empty string if rating is null.\textgreater{}'',\newline
          ``failure\_modes'': [``\textless{}failure\_mode\_1\textgreater{}'', ``\textless{}failure\_mode\_2\textgreater{}'', \ldots{}],\newline
          ``rating'': \textless{}int: 1, 2, or 3, or null if no assistant content\textgreater{}\newline
        \}\}\newline
      ]\newline

      If the turn is rated 3 or null, failure\_modes must be an empty list: [].
\end{tcolorbox}


\subsection{User Speech Fidelity}
\label{app:prompt-user-speech-fidelity}

\begin{tcolorbox}[promptbox, title={\small User Speech Fidelity Prompt}]

        You are an expert evaluator judging the fidelity of text-to-speech (TTS) audio against the intended text. You will listen to one audio clip and verify that the spoken content faithfully reproduces the intended text, with special attention to TTS-critical entities.\newline

        \#\# Evaluation Mode: User\newline

        \#\# Intended Turns\newline
        \{intended\_turns\_formatted\}\newline

        \#\# Understanding the Intended Text\newline

        The intended text may contain non-spoken tags and markers. You must understand these to evaluate fairly.\newline

        \#\#\# Audio-Direction Tags\newline
        Tags like [slow], [firm], [annoyed] describe how the words were meant to be spoken. They are NOT spoken aloud and should never be expected in the audio.\newline

        \#\#\# Interruption Tags\newline
        \{interruption\_tags\_reference\}\newline

        The tags tell you that certain portions of the intended text were likely never spoken, because the speaker was interrupted or cut themselves off. Do NOT penalize for missing words that fall in a region the tags indicate was not spoken.\newline

        \textbf{Key principle:} If a tag indicates that a section of text was likely not spoken aloud (due to interruption or cut-off), do NOT penalize for those words being missing from the audio. Only evaluate fidelity for words that were reasonably expected to have been spoken.\newline

        \#\# Evaluation Criteria\newline

        \#\#\# TTS-Critical Entities (check these carefully)
        \begin{itemize}
          \item \textbf{Personal names}: ``John Smith'' vs ``Jim Smith''
          \item \textbf{Dates and times}: ``December 15th'' vs ``December 50th'', ``3:45 PM'' vs ``3:15 PM''
          \item \textbf{Reference codes}: Confirmation numbers, incident numbers, booking IDs (e.g., ``QWMN62'' vs ``QWN62'')
          \item \textbf{Numeric values}: Dollar amounts, quantities, percentages (e.g., ``\$150'' vs ``\$115'')
          \item \textbf{Addresses}: Street numbers, street names, cities (e.g., ``123 Main Street'' vs ``124 Main Street'')
          \item \textbf{Contact information}: Phone numbers, email addresses (e.g., ``tom\_cobb@gmail.com'')
          \item \textbf{Flight/route numbers}: ``UA204'' vs ``UA240''
          \item \textbf{Serial numbers and other identifiers}
        \end{itemize}

        \#\#\# Error Types
        \begin{itemize}
          \item \textbf{Missing words}: Words in the intended text that were not spoken AND were reasonably expected to have been spoken (i.e., not in a region flagged by interruption tags)
          \item \textbf{Added words}: Extra words spoken that are not in the intended text
          \item \textbf{Wrong words}: Words spoken incorrectly or substituted with different words
          \item \textbf{Entity errors}: Any of the TTS-critical entities above spoken incorrectly
        \end{itemize}

        \#\#\# What to Ignore
        \begin{itemize}
          \item Non-spoken tags: [slow], [firm], [annoyed], and all interruption tags listed above
          \item Words in regions flagged by interruption tags as likely not spoken
          \item Minor pronunciation variations that do not change meaning (accent differences)
          \item Natural filler words (um, uh) if they do not affect core content
          \item Missing words at the END of the LAST turn only (audio recordings are often cut off before the final utterance completes). However, missing words in the middle of the last turn, or missing words in any earlier turn, should still be penalized.
        \end{itemize}

        \#\# Rating Scale (per turn)
        \begin{itemize}
          \item \textbf{3 (High Fidelity)}:
          \begin{itemize}
            \item All expected entities spoken correctly (names, dates, destinations, codes, etc)
            \item All words reasonably expected to have been spoken are present and accurate.
            \item Minor pronunciation variations acceptable.
            \item No audio tags spoken out loud.
          \end{itemize}
          \item \textbf{2 (Medium Fidelity)}:
          \begin{itemize}
            \item All entities spoken correctly (names, dates, destinations, codes, etc)
            \item Part of a turn may be missing (often in the first turn, the first few words are missing)
            \item Some words that were reasonably expected may be missing or spoken slightly incorrectly, but they are not critical and the conversation is able to progress even with this issue.
            \item Potential issues with audio tags being said out loud
          \end{itemize}
          \item \textbf{1 (Low Fidelity)}:
          \begin{itemize}
            \item One or more entity errors (missing entities, incorrect entities, etc) OR
            \item Some other major error that prevents the conversation from continuing in a sensible manner.
          \end{itemize}
        \end{itemize}

        \#\# Response Format\newline
        Respond with a JSON object. Each turn entry must include the turn\_id matching the turn number shown in the Intended Turns above:\newline
        \{\{\newline
          ``turns'': [\newline
            \{\{\newline
              ``turn\_id'': \textless{}int: the turn number from the Intended Turns\textgreater{},\newline
              ``transcript'': \textless{}string: your transcription of the audio for this turn, use only the audio for this not the intended turns\textgreater{},\newline
              ``explanation'': ``\textless{}succinct analysis; for score 1 or 2, quote the specific issue with intended vs actual; note any regions skipped due to interruption tags\textgreater{}'',\newline
              ``rating'': \textless{}1, 2, 3\textgreater{}\newline
            \}\}\newline
          ]\newline
        \}\}
\end{tcolorbox}

\subsection{User Behavioral Fidelity}
\label{app:prompt-user-behavior}

\begin{tcolorbox}[promptbox, title={\small User Behavioral Fidelity Prompt}]

      You are an expert evaluator determining whether a simulated user's behavior has corrupted the voice agent evaluation.\newline

      Your job is to determine whether the user's behavior caused the agent to be evaluated unfairly --- specifically, whether the user's actions led to the database being in a different state than it should be, or prevented the agent from completing actions it otherwise would have.\newline

      \#\# Conversation Evidence\newline
      \{conversation\_evidence\}\newline

      \#\# How the conversation ended\newline
      The conversation ended due to: \{conversation\_end\}\newline

      \#\# User Simulator Instructions\newline
      The following is the full system prompt the user simulator was given for this conversation, including the user's persona, goal, decision tree, must-have criteria, and end-of-call rules. This is the source of truth for what the user-sim was told to do. Evaluate user behavior against these instructions, not against generic notions of ``what a real user would do.''\newline
      \`{}\`{}\`{}\newline
      \{user\_simulator\_instructions\}\newline
      \`{}\`{}\`{}\newline

      \#\# Modification Tools\newline
      The following are the tools that modify database state. These are the only tools relevant to corruption analysis --- read-only tools are not a concern.\newline
      \{modification\_tools\}\newline

      \#\# Evaluation Criteria\newline

      \#\#\# Guiding principle: judge the user from the user's perspective\newline
      The user is a simulated caller who only hears the agent's spoken words. They cannot see tool calls, tool responses, or the agent's internal state. Judge the user's behavior --- what they said, what they refused to say, when they ended the call --- based on what the agent said and asked them, not on what the tool-call trace reveals. Use the tool-call trace asymmetrically:
      \begin{itemize}
        \item \textbf{To exonerate the user (always allowed):} confirm that an agent failure occurred that the user couldn't have known about --- e.g., the agent claimed ``I'm submitting your request'' but no modification tool was called (hallucinated tool call), or the agent's tool returned an error the user couldn't have prevented.
        \item \textbf{To convict the user (use carefully):} only when the user's \textit{visible behavior} (their words, their refusals, their choice to end) plausibly caused the modification problem. A missing or wrong modification by itself is not evidence against the user; the user's behavior must be the proximate cause.
      \end{itemize}

      Analyze the conversation for the following corruption scenarios:\newline

      \#\#\# Corruption Type 1: User invented requests that caused extra modifications\newline
        The user made requests OUTSIDE of their assigned goal that caused the agent to call one or more modification tools listed above.
      \begin{itemize}
        \item Only flag this if the user's off-script request directly led to a modification tool being called.
        \item If the user went off-script but the agent only called read-only tools (e.g., searching, looking up information), this is NOT corruption.
      \end{itemize}

      \#\#\# Corruption Type 2: User ended the conversation prematurely\newline
        The user ended the conversation before the agent had the opportunity to complete the necessary modification tools to fulfill the user's goal.\newline
        Applicability:
      \begin{itemize}
        \item This type only applies when the conversation ended due to the user calling the end call tool.
        \item If the conversation ended due to the agent failing to respond to the user's final turn, this corruption type should never be flagged since the user wasn't the one to end the call.
      \end{itemize}

      Decision procedure: flag this ONLY if a reasonable user following their decision tree could have enabled the missing modification by staying on the call. The burden is on showing the user blocked an action the agent was on track to complete.\newline

      Do NOT flag (agent failures that exonerate the user, even if a required modification did not happen):
      \begin{itemize}
        \item \textbf{Agent error or retry offer.} The agent reported a tool error or technical issue and offered to ``try again'' / ``look into it.'' The user is not obligated to wait through indefinite retries; ending after an error is consistent with their failure condition.
        \item \textbf{Hallucinated tool call.} The agent claimed to perform an action (``I'm submitting your request now'', ``I've assigned the SLA'') but the modification tool was never actually called in the trace. Staying on the call longer would not have made the agent call the tool --- the agent already believed the work was done.
        \item \textbf{Off-script demand.} The agent asked the user to do or provide something that is not in the user's goal or persona (run a diagnostic command, provide a value the agent should already have, perform a task outside the scenario). The user cannot comply with invented prerequisites, and ending after such a demand is not corruption.
        \item \textbf{Skipped wrap-up before the agent's own goodbye.} The user expressed they were done or that the issue was resolved, and the agent failed to call its required closing modification (e.g., mark\_resolved, close\_ticket) before the goodbye exchange. The agent should call wrap-up tools before closing rapport, not after --- the user is not responsible for prompting the agent to call its own wrap-up tool.
        \item \textbf{Agent stuck or unhelpful} for multiple consecutive turns, or the agent said it could not help.
        \item \textbf{Resolution condition met.} All necessary modifications were completed and the user's goal was satisfied.
      \end{itemize}

      DO flag:
      \begin{itemize}
        \item The agent was actively and competently progressing (asking for a specific piece of information that exists in the user's goal, about to call a tool) and the user cut it off without cause.
        \item The user's decision tree explicitly required waiting for a confirmation / request ID, and the user ended without one despite the agent making normal progress (no error, no stuckness, no off-script demand, no hallucinated tool call).
        \item Special case --- transfer to a live agent:
        \begin{itemize}
          \item The user is instructed to always end the call AFTER the agent transfers to a live agent. This should not be penalized.
          \item However the user must give the agent a chance to actually call the transfer\_to\_agent tool. If the user says ``Ok go ahead and transfer me. Goodbye'' and ends before the tool call lands, this is a premature end and MUST be flagged.
        \end{itemize}
      \end{itemize}

      \#\#\# Corruption Type 3: User failed to provide required information\newline
        The user failed to provide information from their goal that the agent explicitly asked for, preventing the agent from completing a necessary modification tool call.
      \begin{itemize}
        \item Only flag this if the agent clearly asked for specific information that was available in the user's goal, the user failed to provide it, and this directly prevented a modification tool from being called.
        \item Do NOT flag this if the agent never asked for the information.
      \end{itemize}

      \#\#\# Corruption Type 4: User looping caused duplicate modifications\newline
        The user repeatedly made the same request in a loop, causing the agent to call the same modification tool multiple times when it should have only been called once.
      \begin{itemize}
        \item Only flag this if the looping directly caused duplicate or extra modification tool calls.
        \item If the user looped but the agent handled it correctly (did not call extra modification tools), this is NOT corruption.
      \end{itemize}

      \#\#\# Corruption Type 5: User violated decision tree instructions causing a wrong modification\newline
        The user explicitly violated a specific instruction in their decision tree (negotiation behavior, edge cases, escalation behavior, resolution condition, or failure condition) AND this violation directly caused a modification tool to be called with different parameters than it would have been if the user had followed their instructions correctly.
      \begin{itemize}
        \item Examples: the user accepted an option that did not meet their must-have criteria when they should have rejected it; the user ignored an edge case instruction (e.g., accepted a standby flight when told to reject standby) and this led to a modification; the user failed to follow their failure condition and instead accepted an unsuitable resolution.
        \item Only flag this if the violation directly caused a modification tool to be called incorrectly. If the user deviated from instructions but no modification tool was affected, this is NOT corruption.
        \item Do NOT flag this if the agent only presented options that failed to meet the user's must-have criteria AND the user had no correct option to choose --- in that case the agent failed, not the user. Only flag this if the user had a correct action available (e.g., rejecting all options, asking for alternatives, triggering the failure condition) but chose incorrectly instead.
      \end{itemize}

      \#\# Rating\newline

      \textbf{Binary Rating:}
      \begin{itemize}
        \item \textbf{1 (Clean)}: The user's behavior did not corrupt the agent evaluation. None of the corruption types above occurred. Minor deviations from the user's instructions that did not affect database state are acceptable.
        \item \textbf{0 (Corrupted)}: One or more corruption types occurred --- the user's behavior caused the agent to be evaluated against an incorrect database state.
      \end{itemize}

      Respond in JSON format:\newline
        \{\{\newline
          ``corruption\_analysis'': \{\{\newline
            ``extra\_modifications'': \{\{``analysis'': ``\textless{}reasoning about whether the user made off-script requests that caused modification tool calls\textgreater{}'', ``detected'': \textless{}bool\textgreater{}\}\},\newline
            ``premature\_ending'': \{\{``analysis'': ``\textless{}reasoning about whether the user ended the call before the agent could complete necessary modifications\textgreater{}'', ``detected'': \textless{}bool\textgreater{}\}\},\newline
            ``missing\_information'': \{\{``analysis'': ``\textless{}reasoning about whether the user failed to provide requested information that blocked a modification\textgreater{}'', ``detected'': \textless{}bool\textgreater{}\}\},\newline
            ``duplicate\_modifications'': \{\{``analysis'': ``\textless{}reasoning about whether user looping caused duplicate modification tool calls\textgreater{}'', ``detected'': \textless{}bool\textgreater{}\}\},\newline
            ``decision\_tree\_violation'': \{\{``analysis'': ``\textless{}reasoning about whether the user violated a specific instruction and whether a correct action was available, and whether this caused an incorrect modification\textgreater{}'', ``detected'': \textless{}bool\textgreater{}\}\}\newline
          \}\},\newline
          ``rating'': \textless{}int: 0 or 1\textgreater{}\newline
        \}\}
\end{tcolorbox}

\subsection{Speakability}
\label{app:prompt-speakability}

\begin{tcolorbox}[promptbox, title={\small Speakability Prompt}]

      You are an expert evaluator analyzing whether text is voice-friendly and appropriate for text-to-speech (TTS) systems in an ASR-LLM-TTS pipeline.\newline

      \#\# Your task: identify violations that make text unsuitable for speech output\newline
      The text MUST be scored as 0 (Voice-Unfriendly) if it contains ANY of the following violations:\newline

      \#\#\# VIOLATIONS (Score = 0):

      \begin{enumerate}
        \item \textbf{Markdown / visual formatting}: Any syntax that creates visual structure but cannot be spoken:
        \begin{itemize}
          \item Bold/italic: ``**important**'' or ``*note*''
          \item Headers: ``\#\# Title'' or ``\# Section''
          \item Markdown tables
          \item Repeated Punctuation/Symbols: Strings of characters (e.g., ----- or *****) that are typically used for visual emphasis
        \end{itemize}
        \item Other kinds of formatting that should not be spoken --- JSONs with brackets, etc.
        \item Missing spaces between words that would cause a TTS system to fail (e.g., ``eighttwentypm'' instead of ``eight twenty PM''). Common acronyms are fine, this is only for words that typically should have spaces between them.
        \item Emojis
      \end{enumerate}

      \#\# Instructions\newline

      Carefully review each assistant turn below. Check each turn for any of the above violations.\newline

      \textbf{If you find even any violations in a turn, the rating for that turn MUST be 0.}\newline

      \#\#\# Assistant Turns\newline
      \{assistant\_turns\_formatted\}\newline

      \#\#\# Response Format\newline
      Provide your response as a valid JSON array, one entry per turn. Each entry must include the turn\_id matching the turn number shown above.\newline
      [\newline
        \{\{\newline
          ``turn\_id'': \textless{}int: the turn number\textgreater{},\newline
          ``explanation'': ``\textless{}string: 1--3 sentence analysis of the speakability of the assistant response, citing specific example of any issues that you detect\textgreater{}'',\newline
          ``rating'': \textless{}int: 0 if ANY violation found, 1 if perfectly voice-friendly\textgreater{}\newline
        \}\}\newline
      ]
\end{tcolorbox}

\subsection{Transcription Accuracy (Key Entities)}
\label{app:prompt-transcription-accuracy}

\begin{tcolorbox}[promptbox, title={\small Transcription Accuracy (Key Entities) Prompt}]

      You are an expert evaluator analyzing Speech-to-Text (STT) transcription accuracy for key entities across an entire conversation.\newline

      Your task:
      \begin{enumerate}
        \item For EACH user turn, identify all key entities in the EXPECTED text
        \item Check if each entity appears CORRECTLY in the TRANSCRIBED text
        \item Mark each entity as correct or incorrect
        \item For entities in regions that were likely never spoken aloud (as indicated by interruption tags), still include them in the output but mark them as skipped
      \end{enumerate}

      \#\# What Counts as an Entity\newline
      An entity must have a \textbf{specific, concrete value} --- something that could be passed as an input to a program or tool (not an AI, but a script or database lookup). Ask yourself: could this value be stored in a variable and used programmatically?

      \begin{itemize}
        \item Names (people, places, organizations): e.g.\ ``John Smith'', ``Austin'', ``Delta Airlines''
        \item Specific dates and times: e.g.\ ``December 15th'', ``3:45 PM'' --- NOT vague references like ``tomorrow morning'' or ``later today''
        \item Confirmation codes / reference numbers: e.g.\ ``ABC123'', ``ZK3FFW''
        \item Flight numbers: e.g.\ ``UA 204''
        \item Amounts and prices: use the specific value only, e.g.\ ``\$120'' --- for qualifier phrases like ``under \$120'', only use the specific value
        \item Addresses: e.g.\ ``123 Main Street''
        \item Phone numbers: e.g.\ ``555-867-5309''
        \item Email addresses: e.g.\ ``john@example.com''
        \item Other specific identifiers: seat numbers, loyalty numbers, booking IDs, etc.
      \end{itemize}

      \textbf{Not an entity:} vague temporal words (``tomorrow'', ``next week'', ``morning''), general descriptors (``the cheap flight'', ``a long trip''), or open-ended qualifiers (``less than an hour'', ``around noon'').\newline

      \#\# Understanding Tags in the Expected Text\newline

      The expected text may contain non-spoken tags and markers. These are metadata --- they were never said aloud and must not be treated as entities or evaluated.\newline

      \#\#\# Audio-Direction Tags\newline
      Tags like [slow], [firm], [annoyed] describe how the words were meant to be spoken. Ignore them entirely.\newline

      \#\#\# Interruption Tags\newline
      \{interruption\_tags\_reference\}\newline

      These markers indicate that parts of the expected text may never have been spoken aloud, because the user was interrupted or talked over. An entity that was never spoken cannot be correctly transcribed, so you must NOT penalize for entities in regions that were likely not said, instead mark them as skipped.\newline

      \textbf{Key principle:} Only evaluate entities that were reasonably expected to have been spoken aloud. If a tag indicates the user was interrupted or talked over before or during an entity, still include the entity in your output but set \texttt{skipped: true} and explain why in the analysis. The \texttt{correct} field should reflect your best assessment of whether the transcription matched, but skipped entities will be excluded from accuracy metrics downstream.\newline

      \#\# User Turns to Evaluate\newline
      \{user\_turns\}\newline

      \#\# Correctness Criteria
      \begin{itemize}
        \item Entity must be present (not missing) --- unless in a region flagged by interruption tags
        \item Entity value must match (minor formatting variations OK)
        \item Numbers: ``150'' and ``one hundred fifty'' are equivalent
        \item Dates: ``December 15th'' and ``Dec 15'' are equivalent
        \item Names: Case-insensitive exact match required
      \end{itemize}

      Important note: The expected text will often feature things formatted like ``one two three'' instead of ``123''. Your goal is to evaluate the semantic equivalence, meaning these are considered equivalent if they were heard in audio.\newline

      \#\# Examples\newline

      \textbf{Example Input:}\newline
      Turn 1:\newline
      Expected: \texttt{My confirmation is A B C one two three on December 15th.}\newline
      Transcribed: \texttt{My confirmation is ABC123 on December 15th.}\newline

      Turn 2:\newline
      Expected: \texttt{Transfer one hundred fifty to account 1 2 3 4 5.}\newline
      Transcribed: \texttt{Transfer \$115 to account 12345.}\newline

      Turn 3:\newline
      Expected: \texttt{[slow] The code is X X F six O H, with the letter O, [assistant interrupts] not zero.}\newline
      Transcribed: \texttt{The code is X... X F 6 O H with the letter O.}\newline

      Turn 4:\newline
      Expected: \texttt{My phone number is four zero four five five five [assistant interrupts] zero eight five six.}\newline
      Transcribed: \texttt{My phone number is 404-555.}\newline

      \textbf{Example Response:}\newline
      [\newline
      \{\{\newline
          ``turn\_id'': 1,\newline
          ``entities'': [\newline
          \{\{\newline
              ``type'': ``confirmation\_code'',\newline
              ``value'': ``A B C one two three'',\newline
              ``transcribed\_value'': ``ABC123'',\newline
              ``analysis'': ``Matches exactly'',\newline
              ``correct'': true,\newline
              ``skipped'': false\newline
          \}\},\newline
          \{\{\newline
              ``type'': ``date'',\newline
              ``value'': ``December 15th'',\newline
              ``transcribed\_value'': ``December 15th'',\newline
              ``analysis'': ``Matches exactly'',\newline
              ``correct'': true,\newline
              ``skipped'': false\newline
          \}\}\newline
          ],\newline
          ``summary'': ``All 2 key entities transcribed correctly.''\newline
      \}\},\newline
      \{\{\newline
          ``turn\_id'': 2,\newline
          ``entities'': [\newline
          \{\{\newline
              ``type'': ``amount'',\newline
              ``value'': ``one hundred fifty'',\newline
              ``transcribed\_value'': ``\$115'',\newline
              ``analysis'': ``Amount wrong: \$150 vs \$115'',\newline
              ``correct'': false,\newline
              ``skipped'': false\newline
          \}\},\newline
          \{\{\newline
              ``type'': ``account\_number'',\newline
              ``value'': ``1 2 3 4 5'',\newline
              ``transcribed\_value'': ``12345'',\newline
              ``analysis'': ``Matches exactly'',\newline
              ``correct'': true,\newline
              ``skipped'': false\newline
          \}\}\newline
          ],\newline
          ``summary'': ``1 out of 2 entities correct. Amount error.''\newline
      \}\},\newline
      \{\{\newline
          ``turn\_id'': 3,\newline
          ``entities'': [\newline
          \{\{\newline
              ``type'': ``confirmation\_code'',\newline
              ``value'': ``X X F six O H'',\newline
              ``transcribed\_value'': ``X F 6 O H'',\newline
              ``analysis'': ``Missing one X --- transcribed 5 characters instead of 6. The code appears before the [assistant interrupts] tag so it is evaluated normally.'',\newline
              ``correct'': false,\newline
              ``skipped'': false\newline
          \}\}\newline
          ],\newline
          ``summary'': ``1 entity found before interruption, partially incorrect (missing one X). No entities after [assistant interrupts] tag to skip.''\newline
      \}\},\newline
      \{\{\newline
          ``turn\_id'': 4,\newline
          ``entities'': [\newline
          \{\{\newline
              ``type'': ``phone\_number'',\newline
              ``value'': ``four zero four five five five zero eight five six'',\newline
              ``transcribed\_value'': ``404-555'',\newline
              ``analysis'': ``The full number is 404-555-0856. The [assistant interrupts] tag appears after `five five five', meaning the last four digits (`zero eight five six') were likely drowned out by the agent speaking over the user. The transcription captured the portion before the interruption. Skipping because the entity spans into the interrupted region and cannot be fully evaluated.'',\newline
              ``correct'': false,\newline
              ``skipped'': true\newline
          \}\}\newline
          ],\newline
          ``summary'': ``1 entity found. Phone number spans into interrupted region --- skipped. Partial transcription (404-555) matches the portion before the interruption.''\newline
      \}\}\newline
      ]\newline

      \#\# Response Format\newline
       Respond with a JSON object. Each turn entry must include the turn\_id matching the turn number shown in the User Turns to Evaluate section above:\newline
      [\newline
      \{\{\newline
          ``turn\_id'': \textless{}int: the turn number from the User Turns to Evaluate section\textgreater{},\newline
          ``entities'': [\newline
          \{\{\newline
              ``type'': ``\textless{}name|date|time|confirmation\_code|flight\_number|amount|address|phone|email|etc\ldots{}\textgreater{}'',\newline
              ``value'': ``\textless{}entity value from expected text\textgreater{}'',\newline
              ``transcribed\_value'': ``\textless{}how it appeared or `missing'\textgreater{}'',\newline
              ``analysis'': ``\textless{}brief reason; if skipped, explain why the entity falls in an interrupted region\textgreater{}'',\newline
              ``correct'': \textless{}true|false\textgreater{},\newline
              ``skipped'': \textless{}true|false\textgreater{}\newline
          \}\}\newline
          ],\newline
          ``summary'': ``\textless{}1--2 sentence summary for this turn\textgreater{}''\newline
      \}\}\newline
      ]
\end{tcolorbox}